O resumo em língua estrangeira é um elemento obrigatório e elaborado conforme
a ABNT NBR 6028. É a tradução do resumo na língua vernácula. De acordo com a
ABNT NBR 14724, todo o texto da monografia deve ser digitado com espaçamento
1,5 entre linhas, \textbf{exceto} as citações de mais de três linhas, notas de rodapé,
referências, legendas das ilustrações e das tabelas, natureza (tipo do
trabalho, objetivo, nome da instituição a que é submetido e área de
concentração), que devem ser digitados ou datilografados em espaço simples.
Texto com fonte Times New Roman ou Arial, justificado, tamanho 12, espaçamento
entre linhas 1,5~cm e conter de 150 a 500 palavras e sem citações. Texto com
fonte Times New Roman ou Arial, justificado, tamanho 12, espaçamento entre
linhas 1,5~cm e conter de 150 a 500 palavras e sem citações. Texto com fonte
Times New Roman ou Arial, justificado, tamanho 12, espaçamento entre linhas
1,5~cm e conter de 150 a 500 palavras e sem citações. Texto com fonte Times
New Roman ou Arial, justificado, tamanho 12, espaçamento entre linhas 1,5~cm e
conter de 150 a 500 palavras e sem citações. Texto com fonte Times New Roman
ou Arial, justificado, tamanho 12, espaçamento entre linhas 1,5~cm e conter de
150 a 500 palavras e sem citações.

\keywords{palavra; palavra; palavra.}

\nota{As palavras-chave devem ser grafadas com iniciais em letra minúscula, com
exceção dos substantivos próprios e nomes científicos, vir logo abaixo do
resumo, separadas por ponto e vírgula e finalizadas com ponto. No mínimo de 3
(três) e no máximo 5 (cinco) palavras-chave.}
